\chapter{Boolean Analysis — Circuit Size}
%%% generated by the blueprint pipeline from TCSlib.BooleanAnalysis.RazborovSmolensky.CircuitSize
%%%%%%%%%%%%%%%%%%%%%%%%%%%%%%%%%%%%%%%%%%%%%%%%%%%%%%%%%%%%%%%%%%%%%%%%%%%%%%%
\section{Overview}

This file identifies the layerwise gate count $\mathrm{gateCountBefore}$ accumulated by the
Razborov--Smolensky induction with the total circuit size $F.\mathrm{size}$ of a feed-forward
circuit, by writing both as the sum of the cardinalities of the non-input layers.  The
identification is then used to restate the polynomial-approximation existence theorems of the
previous chapter with their error bounds phrased directly in terms of the number of gates.

%%%%%%%%%%%%%%%%%%%%%%%%%%%%%%%%%%%%%%%%%%%%%%%%%%%%%%%%%%%%%%%%%%%%%%%%%%%%%%%
\section{Declarations}

\begin{definition}[Index of a non-input layer]
\lean{ACP.gateLayerIdx}
\label{ACP.gateLayerIdx}
\leanok
\uses{ACP.FeedForward}
For a feed-forward circuit $F$, a bound $d \le F.\mathrm{depth}$ and an index
$j : \mathrm{Fin}\,d$, the layer index $\mathrm{gateLayerIdx}\,F\,h_d\,j \in
\mathrm{Fin}(F.\mathrm{depth}+1)$ is $j+1$, i.e.\ the $(j+1)$-st node layer of $F$.  It packages the
non-input layers $1, \dots, d$ as indices into $F.\mathrm{nodes}$ so that they can be summed over.
\end{definition}

\begin{lemma}[Partial gate count as a sum of layer sizes]
\lean{ACP.gateCountBefore_eq_sum_cards}
\label{ACP.gateCountBefore_eq_sum_cards}
\leanok
\uses{ACP.FeedForward, ACP.gateCountBefore, ACP.gateCountBefore_succ, ACP.gateLayerIdx}
\difficulty{5}
Let $F$ be a Boolean feed-forward circuit (alphabet $\bbf_2$, input type
$\mathrm{Fin}\,n$) all of whose node layers are finite, and write $\mathrm{depth}(F)$
for its depth. For a natural number $d$ with $d \le \mathrm{depth}(F)$, the number of
non-input gates in the first $d$ layers equals the total number of nodes across layers
$1,\dots,d$:
\[
(\text{non-input gates in layers } 1,\dots,d) \;=\; \sum_{i=1}^{d} \abs{L_i},
\]
where $L_i$ denotes the set of nodes in the $i$-th layer of $F$.
\end{lemma}

\begin{lemma}[Circuit size as a sum of layer cardinalities]
\lean{ACP.size_eq_sum_cards}
\label{ACP.size_eq_sum_cards}
\leanok
\uses{ACP.FeedForward, ACP.FeedForward.size}
\difficulty{2}
Let $F$ be a layered feedforward circuit over the Boolean alphabet $\bbf_2 = \{0,1\}$
with input type $\mathrm{Fin}\,n$, and suppose every node layer of $F$ is finite. Then
the size of $F$ equals the sum of the cardinalities of its non-input layers:
\[
\mathrm{size}(F) \;=\; \sum_{d\,=\,0}^{\mathrm{depth}(F)-1} \abs{\,\text{layer }
(d+1)\,},
\]
that is, the total number of non-input nodes is obtained by adding up the numbers of
nodes in layers $1, \dots, \mathrm{depth}(F)$.
\end{lemma}

\begin{lemma}[Full-depth gate count equals the circuit size]
\lean{ACP.gateCountBefore_depth_eq_size}
\label{ACP.gateCountBefore_depth_eq_size}
\leanok
\uses{ACP.FeedForward, ACP.FeedForward.size, ACP.gateCountBefore, ACP.gateCountBefore_eq_sum_cards, ACP.gateLayerIdx, ACP.size_eq_sum_cards}
\difficulty{4}
Let $F$ be a feedforward circuit over the alphabet $\bbf_2$ with input type
$\mathrm{Fin}\,n$, and suppose each of its node layers is finite. Counting the non-input
nodes layer by layer through all of the layers $1,\dots,\mathrm{depth}(F)$, and
accumulating these counts over the full depth of $F$, yields exactly the size of $F$,
that is, its total number of non-input nodes.
\end{lemma}

\begin{theorem}[{Polynomial approximation of AC⁰[p] circuits, size form}]
\lean{ACP.exists_poly_distribution_for_circuit_outputs_size}
\label{ACP.exists_poly_distribution_for_circuit_outputs_size}
\leanok
\uses{ACP.ACp_GateOps, ACP.FeedForward, ACP.FeedForward.eval, ACP.FeedForward.onlyUsesGates, ACP.FeedForward.size, ACP.boolInput, ACP.circuitDegreeBound, ACP.exists_poly_distribution_for_circuit_outputs, ACP.gateCountBefore_depth_eq_size}
\difficulty{3}
Let $p$ be a prime and let $F$ be a feedforward circuit over the Boolean alphabet
$\{0,1\}$ with inputs indexed by $\{1,\dots,n\}$ and a finite set of outputs, in which
every layer has only finitely many nodes and every gate belongs to the
$\mathrm{AC}^0[p]$ gate set. Then for every $\ell \in \bbn$ there exist a finite,
nonempty seed set $\mathrm{Seed}$ and, for each seed $s \in \mathrm{Seed}$ and each
output $o$, a polynomial $P_{s,o} \in \bbz/p[X_1,\dots,X_n]$ of total degree at most
$\bigl((p-1)\ell\bigr)^{\mathrm{depth}(F)}$, such that for every Boolean input $x \in
\{0,1\}^n$ the seeds that err on $x$ satisfy
\[
\abs{\bigl\{\, s \in \mathrm{Seed} : P_{s,o}(\bar x) \neq \overline{F(x)(o)}\ \text{for
some output } o \,\bigr\}} \cdot 2^{\ell} \;\le\; \mathrm{size}(F) \cdot
\abs{\mathrm{Seed}},
\]
where $\bar x \in (\bbz/p)^n$ and $\overline{F(x)(o)} \in \bbz/p$ denote the images
under $\{0,1\} \to \bbz/p$ of the Boolean vector $x$ and of the circuit's output at $o$
on input $x$, respectively.
\end{theorem}

\begin{theorem}[{Polynomial approximating distribution for an AC⁰[p] circuit}]
\lean{ACP.exists_poly_distribution_for_circuit_one_size}
\label{ACP.exists_poly_distribution_for_circuit_one_size}
\leanok
\uses{ACP.ACp_GateOps, ACP.FeedForward, ACP.FeedForward.eval₁, ACP.FeedForward.onlyUsesGates, ACP.FeedForward.size, ACP.boolInput, ACP.circuitDegreeBound, ACP.exists_poly_distribution_for_circuit_one, ACP.gateCountBefore_depth_eq_size}
\difficulty{3}
Let $p$ be a prime, and let $F$ be a feedforward circuit over the Boolean alphabet
$\mathrm{Fin}\,2$ with input coordinates indexed by $\mathrm{Fin}\,n$ and a single
output node, in which every layer has finitely many nodes and every gate operation
belongs to the $\mathrm{AC}^0[p]$ gate set (the identity gate, the NOT gate, the
unbounded fan-in AND gate of every arity, and the $\mathrm{MOD}_p$ gate of every arity).
For a Boolean input $x : \mathrm{Fin}\,n \to \mathrm{Fin}\,2$, write $\bar{x} \in
(\bbz/p)^n$ for the point whose $i$-th coordinate is the residue of $x_i$ modulo $p$,
and let $c(x) \in \bbz/p$ be the residue of the single Boolean value that $F$ computes
on $x$. Then for every $\ell \in \bbn$ there exist a finite nonempty index set
$\mathrm{Seed}$ and a family of polynomials $(P_s)_{s \in \mathrm{Seed}}$ in
$\bbz/p[X_1,\dots,X_n]$, each of total degree at most $\bigl((p-1)\ell\bigr)^{d}$, where
$d$ is the depth of $F$, such that for every Boolean input $x$,
\[
\bigl|\{\, s \in \mathrm{Seed} : P_s(\bar{x}) \neq c(x) \,\}\bigr| \cdot 2^{\ell}
\;\le\; \operatorname{size}(F) \cdot \abs{\mathrm{Seed}},
\]
where $\operatorname{size}(F)$ denotes the number of non-input nodes of $F$.
\end{theorem}

\begin{theorem}[{Polynomial approximator list for a single-output $\mathrm{AC}^0[p]$ circuit}]
\lean{ACP.exists_poly_list_for_circuit_one_size}
\label{ACP.exists_poly_list_for_circuit_one_size}
\leanok
\uses{ACP.ACp_GateOps, ACP.FeedForward, ACP.FeedForward.eval₁, ACP.FeedForward.onlyUsesGates, ACP.FeedForward.size, ACP.boolInput, ACP.circuitDegreeBound, ACP.exists_poly_list_for_circuit_one, ACP.gateCountBefore_depth_eq_size}
\difficulty{3}
Let $p$ be a prime and let $F$ be a feedforward circuit over the alphabet
$\mathrm{Fin}\,2$ with $n$ Boolean input variables, finitely many nodes in each layer,
and a single output node, every gate of which belongs to the $\mathrm{AC}^0[p]$ gate
set. Then for each $\ell \in \bbn$ there exists a nonempty list of polynomials $P_1,
\dots, P_m \in (\bbz/p)[X_1, \dots, X_n]$, each of total degree at most
$\bigl((p-1)\ell\bigr)^{\mathrm{depth}(F)}$, with the following property. Given a
Boolean input $x \colon \mathrm{Fin}\,n \to \mathrm{Fin}\,2$, let $\hat{x} \colon
\mathrm{Fin}\,n \to \bbz/p$ denote $x$ cast coordinatewise into $\bbz/p$, and let $b(x)
\in \bbz/p$ denote the image of the single output bit that $F$ computes on $x$. Then for
every such $x$ the number of indices $j$ at which the polynomial disagrees with the
circuit satisfies
\[
\bigl|\{\, j : P_j(\hat{x}) \ne b(x) \,\}\bigr| \cdot 2^{\ell} \;\le\; \mathrm{size}(F)
\cdot m,
\]
where $\mathrm{size}(F)$ is the number of non-input nodes of $F$ and $m$ is the length
of the list.
\end{theorem}
